\documentclass[11pt]{article}
\usepackage[a4paper,margin=25mm]{geometry}
\usepackage{amsmath,amssymb,amsthm,hyperref}
\newtheorem{theorem}{Theorem}
\newtheorem{lemma}[theorem]{Lemma}
\newcommand{\Q}{\mathbb Q}
\newcommand{\Z}{\mathbb Z}
\newcommand{\BB}{\operatorname{BB}}
\newcommand{\Tr}{\operatorname{Tr}}
\title{A fixed 163-variable quartic with a sharper explicit Busy Beaver index}
\author{Z diophantine in Q project}
\date{20 November 2026}
\begin{document}\maketitle
\begin{abstract}
The numerical Busy Beaver specialization established in the supplied
Version 22 manuscript is improved from index \(2^{1500}\) to an explicit
249-digit index less than \(2^{827}\). The coefficient upper bound
improves from \(\BB(2^{2000})\) to \(\BB(2^{1012})\).
The reduction is achieved by evaluating the quantitative Vojta constants
exactly and sharpening the hardwired-record and compiler estimates.
All source arithmetic and geometric hypotheses are retained.
\end{abstract}

\section{Statement and provenance}
Fix the two-symbol, blank-tape, one-tape, two-way-infinite-tape
halting-time Busy Beaver convention with an uncounted halting state,
counted halting transition and permitted unused working states.
Put
\[
t_0=2^{20}(2^{620}+2^{186}+1)(2^{186}+3),\quad
N=\BB(t_0),\quad C=149^N.
\]
Let \(F_N\) be the 163-variable quartic given by the complete coefficient
tables in the project v82 package, including the input as coordinate one.
There are exactly three \(C\)-dependent coefficients:
\[
[x_{78}^2x_{90}]F_N=-6615898C^2,\quad
[x_{78}x_{90}]F_N=-298C^2,\quad
[x_{78}^2]F_N=22201C^4.
\]
All remaining 2786 coefficients are fixed integers.

\begin{theorem}\label{thm:main}
Under the inherited arithmetic, weighted-height geometry, and
register-to-binary compiler lemmas of the uploaded Version~22
manuscript, the fixed polynomial \(F_{\BB(t_0)}\)
defines \(\Z\) Diophantinely in \(\Q\).
It has 163 total variables, degree four, 2789 nonzero monomials and
integer content one. Moreover
\[
2^{826}<t_0<2^{827},\qquad
H(F_{\BB(t_0)})=149^{4\BB(t_0)+2}<\BB(2^{1012}).
\]
\end{theorem}

\paragraph{Research provenance and credit.}
The explicit Busy Beaver record method used here was supplied by
\emph{Elliot Glazer's agentic research notes and Version~22 manuscript}
\cite{v22}. That manuscript identifies the new Busy Beaver estimate
and receiving-source checks as work of an OpenAI Codex agent in
Glazer's research program. The underlying 163-variable scalar-point
quartic and its coefficient tables are due to Jos\'e Brox, building
on Gabe Conant's shared-product constructions with Astra~6.2.
The new evaluations of constants and machine bounds in this note
are subsequent refinements. These are distinct contributions.

The supplied source is the pinned
\texttt{quartic\_v22.tex}, original commit
\texttt{4b43ba288b9b8639cc5cc85f8c1507d64c698659};
its complete SHA-256 is recorded in the accompanying project report.
Its quantitative record construction is in its
Version~22 Part~I, and its original compiler precedes the
physical tables. The inherited project proofs are
\texttt{scalar\_point163\_179\_20261120\_v82.tex}
and the corresponding original physical coefficient tables.
No claim is made that the current computation alone has
independently reproved the complete inherited arithmetic foundations.

\section{Quantitative geometry}
Write \(E:y^2=x^3-(3-\sqrt2)^2x\) over
\(K=\Q(\sqrt2)\), \(R_0=(1+\sqrt2,2\sqrt2)\)
and \(P_0=2R_0\).
For four finite points \(Q_i=[m_i]R_0+\tau_i\),
where \(\tau_i\in E[2](K)\), put
\[
S=\sum_{i=1}^4m_i^2,\quad
T=7+\tfrac14\Tr\bigl(x(Q_1)+r x(Q_2)+x(Q_3)+x(Q_4)\bigr),
\qquad r\in\Q_{>0}.
\]
Freezing zero-index finite torsion values yields at most 256
strata. With \(k\ge1\) active coordinates the graph
\(V\subset E_0^{2k}\times(\mathbf P^1_r\times\mathbf P^1_T)\)
satisfies
\[
\dim V=2k+1\le9,\quad m=\dim V+1\le10,\quad
N_A+1\le3^8,\quad N_P+1=4,
\quad \deg V\le10!3^8.
\]
The source establishes \(h(V)<2^{150}\),
\(h_{\rm add},C_{\rm haut}<100\), and
\(\widehat h_{A_k}=\alpha S\), \(\alpha\ge3/4\).
Its positive-weight pole argument excludes positive-dimensional
constant translates through every actual rational parameter \(r>0\).
The graph and height hypotheses of von Buhren's
\emph{quantitative} Theorem~3.1, printed p.~312, therefore apply.

At \(k=4\), let
\[
\begin{split}
\Lambda&=
 2^{m^2}(8m\dim V)^{m\dim V}
 (N_A+1)(N_P+1)(\deg V)^{3m},\\
\psi&=\prod_{j=2}^{m\dim V}(3j+1).
\end{split}
\]
Exact integer arithmetic gives
\[
\operatorname{bitlength}(\Lambda)=2004,\qquad
\operatorname{bitlength}(\psi)=602.
\]
The published formulae are
\[
e_1=8\Lambda^\psi,\quad e_2=2\Lambda^{2\psi},
\quad\mu=4\Lambda^\psi,\quad
e_3=(4(N_A+1))^{2m^2\dim V}
 \Lambda^{8\psi}\max(h(V),h_{\rm add},C_{\rm haut}).
\]
Hence the following explicit upper bounds hold:
\[
\begin{aligned}
\log_2e_1&<3+2004\psi,&
\log_2(e_2^2)&<2+8016\psi,\\
\log_2\mu&<2+2004\psi,&
\log_2e_3&<27150+16032\psi.
\end{aligned}
\]
Since \(\psi<2^{602}\) and \(16032<2^{14}\),
the last bound is strictly less than \(2^{616}-8\).
The other bounds have a larger margin.
The cells of diameter at most \(1/e_1\)
covering the normalized direction cube \([-1,1]^4\)
number at most \(\lceil4e_1+1\rceil^4<(8e_1)^4\).
The quantitative theorem allows at most nine rapidly growing
records in one cell and stratum, so
\[
\log_2\bigl(256\cdot9\lceil4e_1+1\rceil^4\bigr)
<36+8016\psi<2^{616}.
\]

\begin{lemma}\label{lem:U}
For \(U=2^{2^{616}}\), the explicit Vojta constants
\(e_2^2,e_3,\mu\) are less than \(U/100\).
There are fewer than \(U\) records across all strata.
\end{lemma}
\begin{proof}
The preceding bounds and the source's positive-weight
noncancellation lemma prove the statement by the
pigeonhole argument of quantitative Theorem~3.1.
\end{proof}

\section{Discrete record height and the arithmetic threshold}
Enumerate each rational \(r>0\) and four finite encoded points
\(4|m_i|+e_i\) in finite shells.
For a reduced rational \(z\), write
\(b(z)=\operatorname{bitlength}(\max(|\operatorname{num}z|,
\operatorname{den}z))\), so \(h(z)\le b(z)\le2h(z)+1\).
Initialize \(R_*=U\).
Declare a record when
\[
S>UR_*,\qquad S>U(b(r)+b(T)),
\]
then replace \(R_*\) by \(S\).
The record count \(j_*<U\) by Lemma~\ref{lem:U}.
Let \(R_f\) be its last \(S\), or \(U\) if none occurs.
The running-record argument gives
\[
S\le2U(h(r)+h(T))+U(R_f+2).
\]
The actual source labels satisfy \(B^2\le9S+1\).
Therefore
\[
B^2\le A(h(r)+h(T))+J,\qquad
A=U^2,\quad J=U^2(R_f+3).
\]

\begin{lemma}\label{lem:threshold}
The sharpened expression
\[
G(R_f)=U^3(R_f+3)
\]
is a sufficient defining threshold for every integer
\(N\ge G(R_f)\). It improves \(U^{1000}(R_f+3)\)
from the supplied Version~22 proof.
\end{lemma}
\begin{proof}
The unchanged center-one contact bounds give
\(H_2<2^{4000000}<U\) and \(K=2^{330770}<U\).
Put \(c=|a|\), \(\tau=h(\beta)\), and
\(L_N=\lfloor9(2N-1)/4\rfloor+1\ge N\).
For each rational zero,
\[
B\ge c\ge149^{L_N},\quad
\tau\le K149^{-N}c(1+\log c/3),\quad
h(R_s)\le4\tau+(2N+1)\log149+\log1801500.
\]
If \(N\ge U^3(R_f+3)\), the scalar-height term satisfies
\[
\frac{Ah(R_s)}{B^2}
\le45U^3N149^{-2N}
\le45N^2 149^{-2N}
<\tfrac14.
\]
Also \(J+A\log2<U^2(R_f+4)<N<B^2/4\).

If the comparison residual \(\Delta\ne0\),
the original cubic divisibility \(M_j^3\mid a\Delta\)
and \(|a\Delta|\le14B^5\) imply
\[
B^2/2<10U^3B^{35/18},
\qquad\text{hence }B^{1/18}<20U^3.
\]
But \(B\ge149^N\ge149^{U^3}>U^{1000}\)
contradicts this for \(U>2^{16}\).
Thus \(\Delta=0\).
If the input error numerator \(J_X\ne0\), the original
linear divisibility, Lucas bounds, and the same absorption give
\[
\log B\le\tfrac75\tau+10\log U
<0.11c+\tfrac{c\log c}{30}
<0.1c\log c.
\]
The retained Lucas lower bound
\(\log B\ge(c-6)\log(c-3)>0.45c\log c\)
contradicts this. Thus \(J_X=0\), and
the original half-integrality and squarefree denominator argument
gives \(X\in\Z\).
Integer completeness is uniform already for \(N\ge1\)
in the inherited ordered construction.
\end{proof}

\section{Explicit state and coefficient bounds}
\begin{proof}[Proof of Theorem~\ref{thm:main}]
The Version~22 reference selector has 1506 AST nodes, 16 functions,
call depth 10, no imports or dynamic calls and maximum tuple length 3.
Change its record scale from \(2^{2^{1000}}\) to \(2^{2^{616}}\)
and the arithmetic threshold exponent from 1000 to 3.
The source's register-machine construction has fewer than
\(2^{186}\) fixed core instructions and counters.
Hardwiring \(0\le j_*<U\) uses at most \(2^{616}\) bits;
the two-scratch-register loader uses fewer than 16 counter
instructions per bit. Thus
\[
L<2^{620}+2^{186},\qquad R<2^{186}+2.
\]
The reviewed unary-block/binary-tape conversion uses fewer than
\(2^{20}(L+1)(R+1)\) working states. Consequently its state count
is strictly less than
\[
2^{20}(2^{620}+2^{186}+1)(2^{186}+3)=t_0<2^{827}.
\]
For the actual last-record ordinal \(j_*\), the resulting program
halts after more than \(G(R_f)\) transitions. Padding its unused
states proves \(G(R_f)<\BB(t_0)\).
Lemma~\ref{lem:threshold} proves the defining assertion.

For the coefficient bound, \(2^{826}<t_0<2^{827}\).
A maximizing \(t_0\)-state halting table uses next-state width
\(827\), so its hardwired \(2t_0\) transition entries, each of
\(830\) bits, occupy fewer than \(2^{838}\) bits.
The source's explicit interpreter and two-counter bit loader
therefore give \(L<2^{843}\), \(R<2^{149}\), and fewer than
\[
2^{20}(L+1)(R+1)<2^{1012}
\]
binary working states. This program computes
\(149^{4\BB(t_0)+2}+1\), clocks more than this many transitions,
and then halts. Padding proves
\[
149^{4\BB(t_0)+2}+1<\BB(2^{1012}).
\]
Since \(\BB(t_0)\ge t_0>53\), the v82 exact height formula
\(\max\{126616122944180358034867798936,149^{4N+2}\}\)
equals \(149^{4N+2}\). The polynomial tables prove
the remaining coefficient and variable assertions.
\end{proof}

\paragraph{Verification scope.}
The accompanying independent gate recomputes the integer
\(\Lambda,\psi\), the record and machine bounds, and the source
selector's AST without executing it. It does not independently
reprove all analytic, arithmetic, or finite-state compiler
lemmas inherited from the 3-MB Version~22 manuscript.

\begin{thebibliography}{9}
\bibitem{v22} Elliot Glazer and attributed project contributors (agentic notes with OpenAI Codex), \emph{Z is Diophantine in Q}, Version 22,
supplied \texttt{quartic\_v22.tex}, commit
\texttt{4b43ba288b9b8639cc5cc85f8c1507d64c698659}.
\bibitem{v82} Z diophantine in Q project,
\emph{A scalar-point sharing construction in 163 and 179 variables},
v82, exact physical and coefficient tables.
\bibitem{VB} J. von Buhren,
\emph{Borne de hauteur semi-effective pour le probl\`eme de
Mordell--Lang dans une vari\'et\'e ab\'elienne},
J. Th\'eorie des Nombres de Bordeaux \textbf{29} (2017),
307--320, Theorem~3.1, printed p.~312.
\end{thebibliography}
\end{document}
